% Source: https://artofproblemsolving.com/community/c3h1343990_geometry
% Stable identity: aops-topic-1343990-q01
Since $AC=CD$, $C$ is the midpoint of the hypotenuse $AD$ of the right triangle $ABD$. The midpoint of the hypotenuse is the centre of the circumcircle (the hypotenuse is a diameter, by Thales' theorem), so it is equidistant from all three vertices:
\[
BC=AC=CD .
\]
(This answers why $AB=BC=AC=CD$: the first equality is given, the rest come from the midpoint of the hypotenuse.)

\begin{asy}
usepackage("amsmath");
size(5.5cm);
defaultpen(fontsize(10pt));
path ra(pair P,pair u,pair v,real s=0.2){return P+s*unit(u)--P+s*unit(u)+s*unit(v)--P+s*unit(v);}
pair B=(0,0), A=(0,2), D=(2*sqrt(3),0), C=(A+D)/2;
draw(A--B--D--cycle,linewidth(0.9));
draw(B--C,red);
draw(ra(B,D-B,A-B,0.2));
draw(arc(A,0.4,-90,degrees(D-A)),blue); label("$60^\circ$",A+0.65*dir(-60),blue+fontsize(9pt));
real t=0.07;
pair m1=(A+C)/2, m2=(C+D)/2, m3=(A+B)/2, m4=(B+C)/2;
draw(m1+t*dir(degrees(D-A)+90)--m1-t*dir(degrees(D-A)+90)); draw(m2+t*dir(degrees(D-A)+90)--m2-t*dir(degrees(D-A)+90));
draw(m3+(-t,0.03)--m3+(t,0.03)); draw(m3+(-t,-0.03)--m3+(t,-0.03));
pair nn=dir(degrees(C-B)+90); draw(m4+t*nn+0.03*unit(C-B)--m4-t*nn+0.03*unit(C-B)); draw(m4+t*nn-0.03*unit(C-B)--m4-t*nn-0.03*unit(C-B));
dot(C);
label("$A$",A,N); label("$B$",B,SW); label("$C$",C,NE); label("$D$",D,SE);
label("\mbox{}",(min(currentpicture,user=true)+max(currentpicture,user=true))/2,black); // keeps the text after the figure black when LaTeX typesets the labels inline
\end{asy}

Together with $AB=BC$ we get $AB=BC=CA$, so triangle $ABC$ is equilateral and
\[
\angle DAB=\angle CAB=\boxed{60^\circ}.
\]
